\chapter{Фиксация меры на носителях: источник, смысл и топологические веса}

Проверяется, выводит ли исправленная программа TOE 6.5 одну родительскую
меру, которая до феноменологии определяет нормировку состояния, априорное
распределение носителей, смысл функционала, веса топологий, статистику полей
и совместный гессиан.

\section{Повторное чтение первичного источника}

TOE 6.5 является программой доказательства, а не завершённой мерой. В нём
предлагается вариация спектральной плотности и свободная энергия
\begin{equation}
 F=S_{\text{эфф}}-T_{\text{эфф}}S_{\text{инф}},
\end{equation}
но \(T_{\rm eff}\), профиль обрезания, масштаб регуляризации, ограничение
мод, пространство плотностей, ядро флуктуаций и фон оставлены задачами для
дальнейшего исследования. Мера на классе сравниваемых многообразий не задана.

\section{Исправленный логарифмический словарь}

Для положительного корреляционного оператора при фиксированном носителе
\begin{equation}
 H_C=-\tau^{-1}\log\widehat C,
 \qquad
 \rho_C=\frac{e^{-\tau H_C}}{\Tr e^{-\tau H_C}}
\end{equation}
точно задаёт каноническое состояние. Это устанавливает нормировку состояния,
но не меру на самих носителях.

\section{Два несовместимых смысла функционала}

Внешний логарифм статистической суммы Гиббса и положительный исходный
тепловой след могут по-разному упорядочивать носители. Один функционал
предпочитает увеличение статистической суммы, другой --- уменьшение
положительного следа. Их нельзя одновременно использовать без нового
относительного веса.

\begin{theorem}[Развилка смыслов]
Логарифмический словарь однозначно определяет состояние на фиксированном
носителе, но не выбирает между термодинамическим функционалом Гиббса и
положительным исходным спектральным действием при вариации самого носителя.
\end{theorem}

\section{Конечные топологические веса}

Сравнение разных топологий зависит от конечных коэффициентов при членах
Эйнштейна, квадрата тензора Вейля и Эйлера. Допустимые конечные контрчлены
могут изменить порядок любых двух однопетлевых определителей. Следовательно,
один определитель не задаёт универсального предпочтения топологии без
заранее выведенных правил перенормировки.

\begin{nogocriterion}[Запрет подмены функционала]
Нельзя объявлять функционал Гиббса и положительный исходный спектральный
след двумя формами одной меры, если их относительный знак и конечные
кривизные коэффициенты не следуют из общего источника.
\end{nogocriterion}

\section{Реестр требований}

Из десяти существенных требований проходят только нормировка состояния на
фиксированном носителе и логарифмический операторный словарь. Не выведены:
априорная мера на метриках и топологиях; единственный смысл функционала;
\(T_{\rm eff}\), \(\tau\) и \(\Lambda\); конечные кривизные связи;
статистика полной полевой и BV-меры; векторное дополнение; общий
кинетический оператор и совместный физический гессиан.

\begin{theorem}[Результат фиксации меры]
Текущий итог равен
\begin{equation}
 \text{нормировка состояния: выполнена},\qquad
 \text{родительская мера: не получена}.
\end{equation}
Архитектура носителя закрывается до новых спектральных сумм и утверждений о
предпочтительной топологии.
\end{theorem}

Тождества Гиббса--Фишера и точные тепловые следы остаются корректными
математическими модулями, но не образуют общей меры.

\section{Условия возможного повторного открытия}

Новая архитектура обязана вывести априорную меру на метриках и топологиях,
параметры температуры и обрезания, все конечные кривизные связи, статистику
полей и BV, векторное дополнение и единственную иерархию функционалов без
свободного относительного веса.

Следующая структурная проверка ---
\path{version5_foundational_relative_architecture_gate}; после неё
граничный источник сохраняется как независимый контроль.

Машинный сертификат сохранён в
\path{s2t_v5_carrier_measure_freeze_gate_results.json}.